\documentclass[11pt,a4paper,oneside]{amsbook}
\usepackage[T1]{fontenc}
\usepackage[utf8]{inputenc}
\usepackage{lmodern}
\usepackage[american]{babel}
\usepackage{amsmath,amssymb,amsthm,mathtools}
\usepackage[all,cmtip]{xy}
\usepackage{microtype}
\usepackage{enumitem}
\usepackage{xspace}
\usepackage{geometry}
\geometry{margin=1.1in}
\usepackage[hidelinks]{hyperref}

% SGA notation used in the translated body fragments.
\DeclareMathOperator{\Hom}{Hom}
\DeclareMathOperator{\Spec}{Spec}
\DeclareMathOperator{\SheafHom}{\mathbf{Hom}}
\DeclareMathOperator{\SheafAut}{\mathbf{Aut}}
\DeclareMathOperator{\pr}{pr}
\DeclareMathOperator{\gr}{gr}
\newcommand{\id}{\mathrm{id}}
\renewcommand{\H}{\mathrm{H}}
\newcommand{\R}{\mathbf{R}}
\newcommand{\PP}{\mathbf{P}}
\newcommand{\CC}{\mathbf{C}}
\newcommand{\ZZ}{\mathbf{Z}}
\newcommand{\kres}{\mathit{k}}
\newcommand{\isomto}{\overset{\sim}{\longrightarrow}}
\newcommand{\isomfrom}{\overset{\sim}{\longleftarrow}}
\let\cal\mathcal
\let\goth\mathfrak
\def\to{\mathchoice{\longrightarrow}{\rightarrow}{\rightarrow}{\rightarrow}}
\def\mto{\mathchoice{\longmapsto}{\mapsto}{\mapsto}{\mapsto}}
\def\To{\mathchoice{\Longrightarrow}{\Rightarrow}{\Rightarrow}{\Rightarrow}}
\def\cf{{cf.}\kern.3em}
\def\Cf{{Cf.}\kern.3em}
\def\ie{{i.e.}\kern.3em}
\def\resp{resp.\xspace}
\def\ptbl{.\kern.2em }
\def\quoi{}
\newcommand{\No}{no.~}

% References outside this exposé retain their source keys and printed numbers.
\makeatletter
\newcommand{\SourceRef}[2]{\@ifundefined{r@#1}{#2}{\ref{#1}}}
\makeatother

\addto\captionsamerican{\renewcommand{\chaptername}{Expos\'e}}
\renewcommand{\thechapter}{III}
\renewcommand{\thesection}{\arabic{section}}
\numberwithin{equation}{section}
\theoremstyle{plain}
\newtheorem{theorem}{Theorem}[section]
\newtheorem{proposition}[theorem]{Proposition}
\newtheorem{lemma}[theorem]{Lemma}
\newtheorem{corollary}[theorem]{Corollary}
\newtheorem{corollaries}[theorem]{Corollaries}
\theoremstyle{definition}
\newtheorem{definition}[theorem]{Definition}
\theoremstyle{remark}
\newtheorem{remarks}[theorem]{Remarks}
\newtheorem{remark}[theorem]{Remark}
\newtheorem*{remarksstar}{Remarks}
\newtheorem*{remarkstar}{Remark}
\newtheorem*{scholium}{Scholium}
\newenvironment{namedstatement}[1]
  {\par\addvspace{.5\baselineskip}\noindent\textsc{#1.}\ \itshape\ignorespaces}
  {\par\addvspace{.5\baselineskip}}

\title[SGA~1, Expos\'e III (English draft)]{Smooth morphisms: extension properties\\
{\normalsize SGA~1, Expos\'e III\\English draft}}
\author{A.\ Grothendieck}
\date{}
\hypersetup{
  pdftitle={SGA 1, Expose III: Smooth morphisms, extension properties (English draft)},
  pdfauthor={A. Grothendieck},
  pdfsubject={Unofficial English translation of SGA 1, Expose III}
}

\begin{document}
\frontmatter
\maketitle

\chapter*{About this draft}
This is an unofficial English translation of Expos\'e~III,
``Morphismes lisses: propri\'et\'es de prolongement'', from SGA~1.
The source volume is A.~Grothendieck and M.~Raynaud,
\textit{Rev\^etements \'etales et groupe fondamental} (SGA~1),
S\'eminaire de g\'eom\'etrie alg\'ebrique du Bois Marie, 1960--61.
It follows the slightly corrected SMF recomposition,
\href{https://arxiv.org/abs/math/0206203v2}{arXiv:math/0206203v2}.

This draft contains all seven sections of the expos\'e, including proofs,
footnotes, diagrams, and the final application to formal and ordinary
smooth schemes over a complete local ring. Statement numbering and
mathematical notation follow the source. References to material outside
this expos\'e retain their source numbers. Apparent source issues have not
been silently repaired; scholarly proofreading remains outstanding.

The translator's contribution is licensed under
\href{https://creativecommons.org/licenses/by-sa/4.0/}{CC BY-SA 4.0}.
The French original remains copyright of its authors and publishers.

\tableofcontents
\mainmatter
\chapter{Smooth morphisms: extension properties}
\label{III}

\section{Formally smooth homomorphisms}
\label{III.1}

In~\SourceRef{II}{II}, we limited ourselves to the consideration of homomorphisms of finite type, and consequently, in local homomorphisms~$A\to B$ of local rings, to the case where~$B$ is isomorphic to a localization of an $A$-algebra of finite type. This case is insufficient for various applications, notably in formal geometry or analytic geometry. For example, the formal power-series ring~$B=A[[t_1,\dots,t_n]]$ has, from the point of view of formal geometry, the properties of a smooth algebra over~$A$. In analytic geometry, the same is true of the local ring of a point~$(y,z)$ of a product $Y\times\CC^n$ considered as an algebra over the local ring of~$y$; moreover, the completion of this algebra is isomorphic to the algebra of formal power series in~$n$ indeterminates over the completion of the base ring~$\cal{O}_x$. This leads to the following definition.

\begin{definition}
\label{III.1.1}
Let~$u\colon A\to B$ be a local homomorphism of local rings (noetherian, as recalled). Suppose~$\kres(B)$ is finite over~$\kres(A)$. We say that~$u$ is a \emph{formally smooth homomorphism}, or that the algebra~$B$ is \emph{formally smooth over}~$A$, if there exists a local finite~$\overline{A}$-algebra~$A'$, free over~$\overline{A}$, such that the local components of the semi-local ring~$\overline{B}\otimes_{\overline{A}} A'=B'$ are $A'$-isomorphic to algebras of formal power series over~$A'$\footnote{For a more general and more conceptual definition, motivated by~\ref{III.2.1} below, \cf EGA~$\mathrm{0}_{\mathrm{IV}}$~19.3.1.}.
\end{definition}

(We denote by~$\overline{A}$, $\overline{B}$ the completed rings of~$A$, $B$.) Since~$B'$ is finite and free over~$\overline{B}$, it is indeed a semi-local ring, a direct sum of complete local rings, each of which is still a free module over~$\overline{B}$, hence has the same dimension as~$\overline{B}$, and therefore as~$B$. It follows that the number of variables~$t_i$ in the formal power-series rings considered in~\ref{III.1.1} is equal to $\dim \overline{B}-\dim \overline{A}=\dim B-\dim A$, and in particular is independent of the local component chosen. One sees at once that it is also the dimension of the ring $B\otimes k=B/\goth{m} B$, where~$k=A/\goth{m}$ is the residue field of~$A$; we shall call it the \emph{relative dimension of}~$B$ \emph{with respect to}~$A$.

\begin{remarks}
\label{III.1.2}
It is evident that Definition~\ref{III.1.1} depends only on the homomorphism on the completions~$\overline{A}\to\overline{B}$ induced by~$A\to B$, which justifies the terminology to some extent. We repent here of Definitions~I~\SourceRef{I.3.2}{3.2}~b) and I~\SourceRef{I.4.1}{4.1}~b), which risk being misleading, and prefer to say ``formally unramified'' respectively ``formally étale'' in the cases considered in these definitions, reserving the terminology ``unramified'' and ``étale'' for the case where~$B$ is a localization of an $A$-algebra of finite type\footnote{Or better, ``essentially unramified'' respectively ``essentially étale'', compare EGA IV 18.6.1.}. The reader will immediately verify that ``formally étale'' is equivalent to ``formally smooth and quasi-finite''. Finally, let us point out that there is a reasonable definition of ``formally smooth'' without any prior hypothesis on the residual extension
$\kres(B)/\kres(A)$ (supposed finite here), encompassing among others the local homomorphisms~$A\to B$ such that~$B$ is \emph{flat} over~$A$ and $B/\goth{m}B$ is a \emph{separable extension} of~$A/\goth{m}=k$ (not necessarily of finite type); for example, a Cohen $p$-ring is formally smooth over the ring of $p$-adic integers. It is the lifting property for homomorphisms (compare~\ref{III.2.1}) that should become the definition in this general case. For the applications we have in view, the case treated in Definition~\ref{III.1.1} will suffice; in what follows, by ``formally smooth'' we shall understand ``with finite residual extension''.
\end{remarks}

\begin{lemma}
\label{III.1.3}
If $B$ is formally smooth over~$A$, then $B$ is flat over~$A$.
\end{lemma}

Since flatness is invariant under completion, we may suppose that~$A$ and~$B$ are complete. Since flatness is invariant under a local flat (hence faithfully flat) extension of the base ring, Definition~\ref{III.1.1} reduces us to the case where~$B$ is a formal power-series algebra over~$A$. But then, as an $A$-module, $B$ is isomorphic to a product of $A$-modules isomorphic to~$A$, hence (the base ring~$A$ being noetherian) is $A$-flat as a product of flat $A$-modules.

\medskip
Let us place ourselves under the conditions of~\ref{III.1.1}. Since the residual extensions of the local components of~$B'$ over~$A'$ are trivial, it follows that $L\otimes_k k'$ is an artinian $k'$-algebra whose local components have trivial residual extensions (where~$L$, $k$, $k'$ are the residue fields of~$A$, $B$, $A'$). This necessary condition for the finite free extension~$A'$ to satisfy the condition stated in~\ref{III.1.1} is also sufficient, as follows at once from~\ref{III.1.4}~(i) and~\ref{III.1.5} below.

\begin{proposition}
\label{III.1.4}
Let $A\to B$ be a local homomorphism of local rings, with finite residual extension, and let $A'$ be a finite local $A$-algebra over~$A$, so that $B'=B\otimes_A A'$ is finite over~$B$, hence a semi-local ring, also noetherian. \textup{(i)} If $B$ is formally smooth over~$A$, then the localizations of~$B$ at its maximal ideals are formally smooth over~$A'$. \textup{(ii)} The converse is true if~$A'$ is free over~$A$.
\end{proposition}

We are immediately reduced to the case where~$A$, $B$ are complete.

(i) Let~$A''$ be a finite free local extension of~$A$ such that the local components of~$B''=B\otimes_A A''$ are formal power-series algebras over~$A''$. Extending scalars $A''\to A''\otimes_A A'\to A'''$, where~$A'''$ is a local component of~$A''\otimes_A A'$, one sees that the local components of~$B''\otimes_{A''} A'''=B\otimes_A A'''$ are formal power-series algebras over~$A'''$. But one also has
\[
B\otimes_A A'''=(B\otimes_A A')\otimes_{A'} A'''=B'\otimes_{A'} A''',
\]
while, since~$A''$ is free over~$A$, $A''\otimes_A A'$ is free over~$A'$, and consequently so is~$A'''$, which is a direct factor of it. This proves that~$B'$ is formally smooth over~$A'$.

(ii) Let~$A''$ be a finite free local $A'$-algebra such that the local components of~$B'\otimes_{A'} A''=B\otimes_A A''$ are formal power-series algebras over~$A''$. Since~$A'$ is free over~$A$, so is~$A''$, and hence~$B$ is formally \emph{smooth} over~$A$.

\begin{proposition}
\label{III.1.5}
Let $A\to B$ be a local homomorphism of local rings, with \emph{trivial} residual extension. For~$B$ to be formally smooth over~$A$, it is necessary and sufficient that~$\overline{B}$ be isomorphic to a formal power-series algebra over~$\overline{A}$.
\end{proposition}

Only the necessity has to be proved, and we may suppose that~$A$ and~$B$ are complete. Let~$\goth{m}$ ($\goth{n}$) be the maximal ideal of~$A$ ($B$), and let~$t_1,\dots,t_n$ be elements of~$\goth{n}$ defining a basis of the vector space
\[
(\goth{n}/\goth{n}^2)/\operatorname{Im}(\goth{m}/\goth{m}^2)=\goth{n}/(\goth{n}^2+\goth{m}B).
\]
These elements therefore define a homomorphism of local $A$-algebras
\[
B_1=A[[t_1,\dots,t_n]]\to B.
\]
We prove that it is an isomorphism. It suffices to prove that for every power~$\goth{m}^q$ of~$\goth{m}$, one obtains an isomorphism after reducing modulo~$\goth{m}^q$ (since~$B_1$ and~$B$ are the projective limits of the corresponding rings reduced modulo~$\goth{m}^q$, with~$q$ variable). Since~$B$ and~$B_1$ are flat $A$-modules, the graded rings associated with the $\goth{m}$-adic filtration are obtained by tensoring over~$k=A/\goth{m}$ with $\gr(A)$ the rings $B_1/\goth{m}B_1$ respectively $B/\goth{m}B$. We are thus reduced to showing that $B_1/\goth{m}B_1\to B/\goth{m}B$ is an isomorphism. In view of~\ref{III.1.3}, we are thereby reduced to the case where~$A$ is a \emph{field}~$k$. On the other hand, if~$A'$ is a finite free local $A$-algebra such that~$B\otimes_A A'$ is a formal power-series algebra over~$A'$ (note that this algebra is local, since the residual extension of~$B$ over~$A$ is trivial), to prove that~$B_1\to B$ is an isomorphism, it suffices to prove that~$B_1\otimes_A A'\to B\otimes_A A'$ is one. This reduces us to the case where~$B$ is already a formal power-series algebra (this reduction should have been made first, before reducing to the case of a base field). But then~$B$ is a regular local ring with coefficient field~$k$, and it is well known (and immediate by considering the graded rings associated with the $\goth{n}_1$-adic and~$\goth{n}$-adic filtrations on~$B_1$ and~$B$) that~$B_1\to B$ is an isomorphism. This completes the proof.

\begin{corollary}
\label{III.1.6}
If~$B$ is formally smooth over~$A$, then there exists a finite local $A$-algebra~$A'$ such that the local components of
\[
\overline{B}\otimes_{\overline{A}}\overline{A'}=\overline{(B\otimes_A A')}
\]
are isomorphic to formal power-series algebras over~$\overline{A'}$.
\end{corollary}
Indeed, if~$L/k$ is the residual extension of~$B/A$, consider an extension~$k'/k$ such that the residual extensions in the $k'$-algebra~$L\otimes_k k'$ are trivial. Let~$A'$ be an algebra finite and free over~$A$ such that~$A'/\goth{m}A'=k'$ (one knows that such an algebra exists, for example by reducing step by step to the case where~$k'/k$ is monogenic, and then lifting to~$A$ the coefficients of the minimal polynomial of a generator of~$k'$ over~$k$). It is local. Then~$B\otimes_A A'$ has, at its maximal ideals, trivial residual extensions over that~$k'$ of~$A'$, and the conclusion follows with the help of~\ref{III.1.5}.

\begin{corollary}
\label{III.1.7}
Let~$A\to B$ be a local homomorphism of local rings. For~$B$ to be formally smooth over~$A$, it is necessary and sufficient that~$B$ be flat over~$A$ and that~$B/\goth{m}B$ be formally smooth over~$k=A/\goth{m}$.
\end{corollary}
Making a suitable finite free local extension~$A'$ of~$A$ and using~\ref{III.1.4}~(ii), we are reduced to the case where the residual extension of~$B/A$ is trivial. We know moreover by~\ref{III.1.4}~(i) and~\ref{III.1.3} that the stated conditions are necessary. For the sufficiency, it suffices to observe that the proof of~\ref{III.1.5} proves, under the hypotheses made here, that~$B$ is a formal power-series algebra over~$A$ (supposing~$A$ and~$B$ complete, which is permissible).

\begin{remark}
\label{III.1.8}
It would not be difficult to develop, for formally smooth homomorphisms, the analogue of all the properties of smooth morphisms studied in~\SourceRef{II}{II}. For the differential properties, however, this requires a modification of the usual definition of Kähler differentials (\cf I~\SourceRef{I.1}{1}), with completed tensor products replacing ordinary tensor products. We shall content ourselves with evoking these abysses here, what precedes being sufficient for our purpose.

It remains to make the link between formal smoothness and the notion of smoothness developed in~\SourceRef{II}{II} (and of which we have not yet made any use):
\end{remark}

\begin{proposition}
\label{III.1.9}
Let $A\to B$ be a local homomorphism, with~$B$ a localization of an $A$-algebra of finite type. For~$B$ to be smooth over~$A$, it is necessary and sufficient that it be formally smooth over~$A$.
\end{proposition}

Using~\ref{III.1.7} and~\SourceRef{II.2.1}{II.2.1}, we are reduced to the case where~$A$ is a field.

Using~\ref{III.1.4}~(ii) and~\SourceRef{II.4.13}{II.4.13}, a suitable extension~$k'$ of~$k$ reduces us to the case where the residual extension for~$B/k$ is trivial. In view of~\ref{III.1.5} (respectively~\SourceRef{II.5.2}{II.5.2}), $B$ is then smooth over~$k$ (respectively formally \emph{smooth over}~$k$) if and only if~$B$ is a regular local ring (respectively its completion is a formal power-series algebra over~$k$). But it is well known that these two conditions are equivalent (the residual extension being trivial).

\section[The lifting property]{The lifting property characteristic of formally smooth homomorphisms}
\label{III.2}

\begin{theorem}
\label{III.2.1}
Let $A \to B$ be a local homomorphism of local rings defining a finite residual extension. The following conditions are equivalent:
\begin{enumerate}
\item[(i)]
$B$ is formally smooth over~$A$
\item[(ii)] For every local homomorphism $A \to C$, where $C$ is a
local ring \emph{complete}, every ideal $J$ of~$C$ contained in
$\goth{r}(C)$, and every local $A$-homomorphism $B \to C/J$, there
exists an $A$-homomorphism (necessarily local) $B \to C$ which
lifts it.
\item[(iii)] For every $A$-algebra $C$ (not necessarily a
noetherian ring), every nilpotent ideal $J$ of~$C$, and every
continuous $A$-homomorphism $B \to C$ (\ie vanishing on a power of
$\goth{r}(B)$), there exists an $A$-homomorphism $B \to C$
(necessarily continuous as well) which lifts it.
\item[(iv)] The same statement as \textup{(ii)} and \textup{(iii)}, but with $C$
a local artinian ring, finite over~$A$.
\item[(iv~bis)] As in \textup{(iv)}, but with $J$ moreover square-zero.
\end{enumerate}
\end{theorem}

\begin{remark}
For the rest of this exposé, we shall use only the implication
(iv)$\To$(i), or (iv~bis)$\To$(i); the direct implication
(i)$\To$(ii) will be proved by another method in the next no. when
$B$ is a localization of an algebra of finite type over~$A$. Let us
recall that in the ``good'' theory of Cohen theorems\footnote{\Cf EGA
$0_{\textup{IV}}$ 19.3, 19.8}, property (ii) or (iii) becomes the
definition of formally smooth homomorphisms, whereas \ref{III.1.1}
becomes a characteristic property valid only in the case of a finite
residual extension. Care should be taken that of properties (ii) and
(iii), neither is more general than the other. One could give an
equivalent property covering both, by introducing a linearly
topologized ring $C$, \emph{separated} and \emph{complete}, a
\emph{closed topologically nilpotent} ideal of~$C$, and a continuous
homomorphism $A \to C$ (thus making $C$ a topological $A$-algebra);
we shall leave this modification to the reader.
\end{remark}

\subsubsection*{Proof of~\ref{III.2.1}}
We shall prove (i)$\To$(iii)$\To$(ii), then (iv)$\To$(i). Since
(ii)$\To$(iv) is trivial, and the equivalence of (iv) and (iv~bis)
is seen by an immediate induction on the integer $n$ such that
$J^n = 0$, this will finish the proof.

(i)$\To$(iii). An immediate induction reduces us to the case where
$J^2 = 0$. Since $C$ is finite over~$A$, there exists a power
$\goth{m}^q$ of the maximal ideal of~$A$ which annihilates $C$.
Dividing by $\goth{m}^q$, and noting that $B/\goth{m}^qB$ is still
formally smooth over~$A/\goth{m}^q$ by \ref{III.1.4}~(i), we may
suppose that $A$ is artinian. Since $B$ is flat over~$A$ by
\ref{III.1.3}, $B$ \emph{is free over} $A$, since $A$ is artinian.
Thus there exists an \emph{$A$-module homomorphism}
$$
w \colon B \to C
$$
which lifts the given homomorphism $u \colon B \to C/J$. Put
$$
f(x,y) = w(xy) - w(x)w(y) \qquad (x,y \in B)
$$
then $f(x,y) \in J$, and $f$ is therefore an $A$-bilinear map from
$B \times B$ into $J$. For there to exist a lift $v \colon B \to C$
of $u$ which is an algebra homomorphism, it is necessary and
sufficient that there exist an $A$-linear map $g \colon B \to J$ such
that $v = w + g$ is an algebra homomorphism, which is written
\begin{align*}
g(1) &= 1 - w(1) \\
g(x,y) - u(x)g(y) -u(y)g(x) &= -f(x,y) \qquad (\text{for } x,y \in
B)
\end{align*}
This is a system of \emph{linear} equations in
$\Hom_A(B,J)$, with right-hand sides in $J$; hence it has a solution
if and only if the corresponding system in $\Hom_A(B,J)\otimes_A A'$, with
right-hand sides in $J' = J \otimes_A A'$, has a solution, $A'$
denoting a faithfully flat algebra over~$A$. Now let $A'$ be an
algebra finite and free over~$A$, local, such that $B' = B \otimes_A A'$ is
a formal power-series algebra over~$A'$ (N.B. in our proof one may
suppose $A$ and $B$ complete, as is immediately seen). Since $A'$ is
free of finite type over~$A$, we have
$$
\Hom_A(B,J) \otimes_A A' = \Hom_{A'}(B',J')
$$
and one verifies that the system of equations obtained in
$\Hom_{A'}(B',J')$ is the one which determines the homomorphisms of
$A'$-algebras $B' \to C' = C \otimes_A A'$ which lift the
homomorphism $u' \colon B' \to C'/J'$ deduced from~$u$ by extension
of scalars, by ``correcting'' by an $A'$-module homomorphism
$g' \colon B' \to J'$ the $A'$-module homomorphism $w' \colon B' \to C'$
deduced from~$w$ by extension of scalars. (Note that $B$ generates
$B'$ as an $A'$-module.) We are thus reduced to proving (iii) when
$B$ is a \emph{formal power-series algebra} over~$A$,
$B = A[[t_1,\dots,t_n]]$. Let us lift arbitrarily the images in
$C/J$ of the $t_i$ to elements $z_i$ of~$C$. Since the $z_i$ modulo
$J$ are nilpotent ($u \colon B \to C/J$ being continuous), the $z_i$
themselves are nilpotent (since $J$ is nilpotent), and hence the
$z_i$ define a continuous homomorphism of topological $A$-algebras
from $B$ into the discrete ring $C$, evidently lifting $u$, qed.

(iii)$\To$(ii).
Let $\goth{n}$ be the maximal ideal of~$C$, and for every integer
$q > 0$, put
$$
C_q = C/\goth{n}^q \text{\quoi, \quad} J_q = (J +
\goth{n}^q)/\goth{n}^q
$$
Thus $C_q/J_q$ identifies with a quotient algebra of~$C/J$. On the
other hand, the composite homomorphism $u_q \colon B \to C/J \to
C_q/J_q$ is continuous from~$B$ into the discrete ring $C_q/J_q$,
and $J_q$ is a nilpotent ideal in~$C_q$. We then construct, step by
step, $A$-homomorphisms
$$
v_q \colon B \to C_q
$$
such that (a) $v_q$ lifts $u_q$ and (b) $v_q$ lifts $v_{q-1}$. The
possibility of the induction is easily verified, since
$$
u_q \colon B \to C/(J + \goth{n}^q) \text{\quad and \quad} v_{q-1}
\colon B \to C/\goth{n}^{q-1}
$$
define the same homomorphism
$$
B \to C/((J + \goth{n}^q) + \goth{n}^{q-1}) = C/(J + \goth{n}^{q-1}) =
C_{q-1}/J_{q-1}
$$
namely $u_{q-1}$, they define a homomorphism
$$
B \to C/J_q' \text{\quad where \quad} J_q' = (J + \goth{n}^q)\cap
\goth{n}^{q-1} \supset \goth{n}^q
$$
(from which both arise by reduction). We are therefore reduced to
lifting a homomorphism $B \to C/J_q'$ from $B$ into a quotient of
$C_q$ by an ideal $J_q'/\goth{n}^q$ contained in $J_q$, hence
nilpotent, and this is possible by hypothesis (iii).

This done, the $v_q$ define a homomorphism from $B$ into the projective
limit $C$ of the $C_q$. Since $J$ is closed, $J$ is the projective
limit of the $J_q$; hence $v$ lifts $u$, qed.

(iv)$\To$(i). One observes first at once that if (iv) is verified,
(iv) remains verified for the local components of $B \otimes_A A'$ over
$A'$, if $A'$ is a finite local algebra over~$A$. Taking $A'$ free
over~$A$ and such that the residual extensions of $B'$ over~$A'$ are
trivial, we are reduced, by \ref{III.1.4}~(ii), to the case where
the residual extension of $B$ over~$A$ is trivial. We shall then prove
the following slightly more precise result:

\begin{corollary}
\label{III.2.2}
Under the conditions of~\ref{III.2.1}, suppose moreover that the residual
extension of $B$ over~$A$ is trivial. Then the equivalent conditions
of~\ref{III.2.1} are also equivalent to the following two conditions
(supposing in \textup{(v)} that $A$ and $B$ are complete):
\begin{enumerate}
\item[(iv~ter)] As in \textup{(iv)}, but with the local artinian ring $C$ finite over~$A$ restricted to have a trivial residual extension (and moreover, if one wishes, with the ideal $J$ square-zero).
\item[(v)] There exists a local $A$-homomorphism (where $n = \dim{\goth{n}/(\goth{n}^2+\goth{m}B})$)
$$
u \colon B \to B_1 = A[[t_1, \dots, t_n]]
$$
inducing an \emph{isomorphism}
$$
\goth{n}/(\goth{n}^2+\goth{m}B) \isomto
\goth{n}_1/(\goth{n}_1^2+\goth{m}B_1)
$$
where $\goth{n}$ ($\goth{n}_1$) is the maximal ideal of $B$
($B_1$), and $\goth{m}$ is that of~$A$.
\end{enumerate}
\end{corollary}

\subsubsection*{Proof}
Since (iv~bis) evidently implies (iv~ter)—setting aside the joke about
the square-zero ideal—it will suffice to prove (iv~ter)$\To$(v)$\To$(i).

(iv~ter)$\To$(v). Choose a basis $a_1,\dots,a_n$ of
$\goth{n}/(\goth{n}^2+\goth{m}B)$, which therefore defines a
local homomorphism of $A$-algebras
$$
B \to B_1/(\goth{n}_1^2+\goth{m}B_1) = k[t_1, \dots, t_n]/(t_1, \dots,
t_n)^2
$$
which can be lifted step by step, by virtue of (iv~ter), to
homomorphisms of $A$-algebras from $B$ into $B_1/\goth{n}_1^2$,
$B_1/\goth{n}_1^3$, and so on; passing to the projective limit gives
the homomorphism $B \to B_1$ with the desired property.

(v)$\To$(i). As in the commutative diagram
$$
\xymatrix{ \goth{m}/\goth{m}^2 \ar[r] \ar[d] & \goth{n}/\goth{n}^2
\ar[r] \ar[d] & \goth{n}/(\goth{n}^2+\goth{m}B) \ar[r] \ar[d] & 0 \\
\goth{m}/\goth{m}^2 \ar[r] & \goth{n}_1/\goth{n}_1^2 \ar[r] &
\goth{n}_1/(\goth{n}_1^2+\goth{m}B_1) \ar[r] & 0}
$$
the two rows are exact, and the extreme vertical arrows are
surjective; the middle arrow is surjective and it follows (since $B$
is complete) that $B \to B_1$ is \emph{surjective}. Let $x_i$
($1 \leq i \leq n$) be elements of~$B$ lifting the $t_i$. They
therefore define a homomorphism of $A$-algebras $B_1 \to B$, which
will be surjective for the same reason as $u$, and whose composite
with $u$ is the identity by construction. Thus $B_1 \to B$ is also
injective, and is consequently an isomorphism. We obtain

\begin{corollary}
\label{III.2.3}
Under the conditions of~\ref{III.2.2}~\textup{(v)}, $u$ is necessarily
an isomorphism.
\end{corollary}

This finishes the proof that $B$ is formally smooth over~$A$. At the
same time we have recovered~\ref{III.1.5} (though there is little
merit in that).

\section{Local infinitesimal extension of morphisms into a smooth $S$-scheme}
\label{III.3}

\begin{theorem}
\label{III.3.1}
Let $f \colon X \to Y$ be a morphism locally of finite type.
The following conditions are equivalent:
\begin{enumerate}
\item[(i)] $f$ is smooth.
\item[(ii)] For every prescheme $Y'$ over~$Y$, every closed
sub-prescheme $Y'_0$ of~$Y'$ having the same underlying space as
$Y'$, every $Y$-morphism $g_0 \colon Y'_0 \to X$, and every
$z \in Y'_0$, there exists an open neighborhood $U$ of~$z$ in~$Y'$
and an extension $g$ of~$g_0|Y'_0 \cap U$ to a $Y$-morphism $U \to
X$.
\item[(ii~bis)] For $Y'$, $Y'_0$, and $z$ as in \textup{(ii)}, putting
$X' = X \times_Y Y'$, $X'_0 = X \times_Y Y'_0$, every section of
$X'_0$ over~$Y'_0$ extends to a section of~$X'$ over an open
neighborhood $U$ of~$z$.
\item[(iii)] For every $Y$-scheme $Y'$, spectrum of a local artinian
ring finite over some $\mathcal{O}_y$ ($y \in Y$), every nonempty
closed sub-prescheme $Y'_0$ of~$Y'$, and every $Y$-morphism
$g_0 \colon Y'_0 \to X$, there exists a $Y$-morphism $g
\colon Y' \to X$ extending $g_0$.
\item[(iii~bis)] For every $Y'$, $Y'_0$ as in \textup{(iii)}, putting
$X' = X \times_Y Y'$, $X'_0 = X \times_Y Y'_0$, every section of
$X'_0$ over~$Y'_0$ extends to a section of~$X'$ over~$Y'$.
\end{enumerate}
\end{theorem}

\subsubsection*{Proof} The equivalence of (ii) and (ii~bis), on the one hand, and of
(iii) and (iii~bis), on the other, is trivial, as is the implication
(ii)$\To$(iii). It remains to prove (i)$\To$(ii) and (iii)$\To$(i).

(i)$\To$(ii). Let $x = g_0(z)$. Replacing $X$ by a suitable open
neighborhood of~$x$, and $Y'$ by the prescheme induced on the open
inverse image of the latter under~$g_0$, we may suppose that $X$ is
\'etale over~$Y[t_1, \dots, t_n]$. Consider the composite $Y$-morphism
$Y'_0 \to X \to Y[t_1, \dots, t_n]$; it is defined by $n$ sections of
the sheaf $\mathcal{O}_{Y'_0}$, which can therefore be extended in a
neighborhood of~$z$ to sections of~$\mathcal{O}_{Y'}$, so we may suppose
that the morphism in question has been extended to a
$Y$-morphism $Y' \to Y[t_1, \dots, t_n]$. By~$\SourceRef{I.5.6}{I.5.6}$ there
is then a unique $Y$-morphism $g \colon Y' \to X$ lifting the preceding
one, and at the same time extending $g_0$, qed.

(iii)$\To$(i). Since the set of points where $f$ is smooth is open, it
suffices to prove that it contains every $x \in X$ that is
\emph{closed} in its fiber. Let $y = f(x)$. Then $\mathcal{O}_x$ is an
algebra over~$\mathcal{O}_y$, a localization of an algebra of finite
type, with finite residual extension. On the other hand, hypothesis
(iii) implies that every homomorphism from $\mathcal{O}_x$ into an
algebra $A/J$, where $A$ is a local artinian algebra finite over~
$\mathcal{O}_y$, and $J$ is an ideal contained in its radical, lifts
to a homomorphism from $\mathcal{O}_x$ into the algebra $A$ (taking
into account that a morphism from $\Spec(B)$, with $B$ local, into
$X$ is determined bijectively by a local homomorphism from an
$\mathcal{O}_x$ ($x \in X$) into $B$). By~\ref{III.2.1} it follows that
$\mathcal{O}_x$ is formally smooth over~$\mathcal{O}_y$, hence smooth
over~$\mathcal{O}_y$ by~\ref{III.1.9}.

\begin{corollary}
\label{III.3.2}
Let $f\colon X\to Y$ be as in~\ref{III.3.1}. The following conditions
are equivalent:
\begin{enumerate}
\item[(i)] $f$ is \'etale.
\item[(ii)] Condition~\textup{(ii)} of~\ref{III.3.1} holds with
{\emph{uniqueness}} of the extension $g$ of~$g_0$ to~$U$.
\item[(iii)] Condition~\textup{(iii)} of~\ref{III.3.1} holds with
{\emph{uniqueness}} of~$g$.
\end{enumerate}
\end{corollary}

It suffices to note, in the proof of (i)$\To$(ii) above, that one can
have uniqueness (when $Y'_0$ is not identical to~$Y'$ in a neighborhood
of~$z$) only if $n=0$, a condition that is known to be sufficient.

\begin{corollary}
\label{III.3.3}
Let $X$ be a prescheme locally of finite type over a {\emph{complete}}
local ring $A$, $y$ the closed point of~$Y=\Spec(A)$, and $x$ a
point of~$f^{-1}(y)$ {\emph{rational}} over~$\kres(y)$. If $X$ is
{\emph{smooth over~$A$ at~$x$}}, then there exists a section $s$ of
$X$ over~$Y$ ``passing through~$x$'' \ie such that $s(y)=x$.
\end{corollary}

In particular, if $X$ is smooth over~$A$, then the natural map
$$
\Gamma(X/Y)\to\Gamma(X\otimes_A k/k)
$$
from the sections of~$X$ over~$Y$ to the set of points of the fiber
$f^{-1}(y)=X\otimes_A k$ rational over~$k$, is surjective. This fact
was especially well known and used when $A$ is a discrete valuation
ring and $X$ is proper over~$A$ (in fact, projective over~$A$), in
which case the sections of~$X$ over~$Y$ (\ie the ``points of~$X$ with
values in~$A$'') also identify with the rational sections, \ie with the
points of~$X\otimes_A K=X_K$ (which is a proper smooth scheme over~$K$)
with values in~$K$, the field of fractions of~$A$, \ie with the points
of~$X$ rational over~$K$.

\section{Local infinitesimal extension of smooth $S$-schemes}
\label{III.4}

\begin{theorem}
\label{III.4.1}
Let $Y$ be a locally noetherian prescheme, let $Y_0$ be a closed
subprescheme having the same underlying space, let $X_0$ be a smooth
$Y_0$-prescheme, and let $x$ be a point of~$X_0$. Then there exist an
open neighborhood $U_0$ of~$x$, a prescheme $X$ smooth over~$Y$, and a
$Y_0$-isomorphism:
$$
h\colon U_0\isomto X\times_Y Y_0.
$$
Moreover, if $(U'_0,X',h')$ is another solution of this problem, then
``it is isomorphic to the first in a neighborhood of~$x$''.
\end{theorem}

We leave it to the reader to make precise what is meant by this. One
may note that for $U_0$ given, a solution of the stated problem
amounts to giving, on~$U_0$, a sheaf of algebras $\cal{B}$ over
$f_0^{-1}(\cal{O}_Y)$ (where $f_0$ is the underlying continuous map
of the structural morphism $U_0\to Y_0$), equipped with a ring
homomorphism $\cal{B}\to \cal{O}_{U_0}$ compatible with the
homomorphism $f^{-1}(\cal{O}_Y)\to f^{-1}(\cal{O}_{Y_0})$, such that

(a) this homomorphism induces an isomorphism
$$
\cal{B}\otimes_{f^{-1}(\cal{O}_Y)} f^{-1}(\cal{O}_{Y_0})\isomto
\cal{O}_{Y_0};
$$

(b) $U_0$ equipped with~$\cal{B}$ becomes a smooth $Y$-prescheme.

In this way, the precise meaning of the assertion of local uniqueness
becomes particularly evident.

\subsubsection*{Proof} We may already suppose that $X_0$ is
\'etale \hbox{over some $Y_0[t_1,\dots,t_n]=Y'_0$}. But the latter
may be regarded as a closed subprescheme of~$Y'=Y[t_1,\dots,t_n]$,
having the same underlying space. By I~\SourceRef{I.8.3}{8.3}, there
exists a prescheme $X$ \'etale over~$Y'$, and a $Y'_0$-isomorphism
$X\times_{Y'}Y'_0\isomto X'$. Existence is thus obtained. For
uniqueness, one uses property~\ref{III.3.1}~(ii) of smooth morphisms,
taking into account the following lemma.

\begin{lemma}
\label{III.4.2}
Let $Y$ be a prescheme, let $Y_0$ be a closed subprescheme defined by
a locally nilpotent sheaf of ideals $\cal{J}$, let $X$ and $X'$ be
$Y$-preschemes, and let $u\colon X\to X'$ be a $Y$-morphism. Suppose
that $X$ is flat over~$Y$. In order that $u$ be an isomorphism, it is
necessary and sufficient that $u_0\colon X\times_Y Y_0\to
X'\times_Y Y_0$ be so.
\end{lemma}

The proof is easy, by passing to the affine case and looking at the
associated graded rings. One should note moreover that the analogous
statement, obtained by replacing ``isomorphism'' by ``closed
immersion'', is also valid, and without the flatness hypothesis.

\begin{remark}
\label{III.4.3}
It is essential to note that the local extension $X$ obtained in~
\ref{III.4.1}
\emph{is not canonical}, in other words the local isomorphism between
two solutions is not unique, \ie in general there exist nontrivial
$Y$-automorphisms of~$X$ inducing the identity on the closed
subprescheme $X_0=X\times_Y Y_0$. This is why, for the construction
of \emph{global} infinitesimal extensions of smooth preschemes, one
must expect the existence of an obstruction of cohomological nature,
which will be made precise below (\No~\ref{III.6}).
\end{remark}

\section{Global infinitesimal extension of morphisms}
\label{III.5}

Let $T$ be a topological space, let $\cal{G}$ be a sheaf of groups
on~$X$, and let $\cal{P}$ be a sheaf of sets on~$T$ on which
$\cal{G}$ acts (on the right, to fix ideas). One says that $\cal{P}$ is
\emph{formally principal homogeneous} under~$\cal{G}$ if the familiar
homomorphism
$$
\cal{G}\times \cal{P}\to \cal{P}\times \cal{P}
$$
of sheaves of sets, deduced from the operations of~$\cal{G}$ on~$\cal{P}$,
is an {\emph{isomorphism}}. This amounts to saying that for every
$x\in T$, $\cal{P}_x$ is {\emph{empty or a principal homogeneous space}}
under the ordinary group~$\cal{G}_x$, or also that for every open set
$U$ of~$T$, $\cal{P}(U)$ is empty or a principal homogeneous space under
the ordinary group~$\cal{G}(U)$. One says that $\cal{P}$ is a
{\emph{principal homogeneous sheaf}} under~$\cal{G}$ if it is formally
so and if, in addition, the $\cal{P}_x$ are nonempty (in other words, if
{\emph{all}} the $\cal{P}_x$ are principal homogeneous spaces, hence
nonempty, under the~$\cal{G}_x$)\footnote{It seems preferable to adopt the shorter and more expressive term
``{\emph{torsor under}}~$\cal{G}$'', introduced in the thesis of
J.\ptbl\textsc{Giraud}.}. Recall that the set of classes (up to
isomorphism) of principal homogeneous sheaves under~$\cal{G}$ identifies
with the cohomology set $\H^1(T,\cal{G})$, which is also the usual
cohomology group of~$T$ with coefficients in~$\cal{G}$ when $\cal{G}$ is
commutative. Thus, for every principal homogeneous~$\cal{P}$, there is
a characteristic class $c(\cal{P})\in\H^1(T,\cal{G})$, whose triviality
is necessary and sufficient for $\cal{P}$ to be trivial (\ie isomorphic
to~$\cal{G}$, on which $\cal{G}$ acts by right translations), or
 equivalently for $\cal{P}$ to have a section.

\begin{proposition}
\label{III.5.1}
Let $S$ be a prescheme, let $X$ and $Y$ be preschemes over~$S$, and let
$Y_0$ be a closed sub-prescheme of~$Y$ defined by an ideal $\cal{J}$ on~$Y$
{\emph{of square zero}}. Let $g_0$ be an $S$-morphism from~$Y_0$ to~$X$,
and let $\cal{P}(g_0)$ be the sheaf on~$Y$ whose sections on an open set
$U$ are the extensions $g\colon U\to X$ of $g_0|U\cap Y_0$ to an
$S$-morphism~$g$. Then $\cal{P}(g_0)$ is a sheaf which is naturally
{\emph{formally principal homogeneous}} under the commutative sheaf of
groups
$$
\cal{G}=\SheafHom_{\cal{O}_{Y_0}}(g_0^*(\Omega^1_{X/S}),\cal{J}).
$$
\end{proposition}

Put $\cal{P}=\cal{P}(g_0)$. For every open set $U$ of~$Y$ we must
define a map
$$
\cal{P}(U)\times \cal{G}(U)\to \cal{P}(U)
$$
so that (a), for fixed $g\in\cal{P}(U)$, the map $s\mto gs$ from~$\cal{G}(U)$
to~$\cal{P}(U)$ is bijective, (b) $\cal{P}(U)$ becomes a set with group
of operators~$\cal{G}(U)$, and (c) the preceding maps are compatible
with the restriction operators for an open set $V\subset U$. The
verification of (c) is trivial, so for simplicity one may suppose
$U=Y$. The verification of (b) (which is, if one wishes, local in
nature) will be left to the reader; we shall limit ourselves, for a
fixed $g\in\cal{P}(Y)$, to defining a natural bijection from~$\cal{G}(Y)$
onto~$\cal{P}(Y)$. Thus suppose that an $S$-morphism
$g\colon Y\to X$ is already given, and seek a canonical bijection
\begin{equation*}
\label{eq:III.5.1.*}
\tag{$*$}
\Hom_{\cal{O}_{Y_0}}(g_0^*(\Omega^1_{X/S}),\cal{J})\isomto\cal{P}(Y)
\end{equation*}
where $\cal{P}(Y)$ is the set of $S$-morphisms $g'$ from~$Y$ to~$X$
inducing the same morphism $g_0\colon Y_0\to X$ as~$g$.

The datum of such a $g'$ is equivalent to the datum of an $S$-morphism
$h\colon Y\to X\times_S X$ such that ${\pr}_1\circ h=g$, and
$h\circ i=(g_0,g_0)$, where ${\pr}_1\colon X\times_S X\to X$ is the
first projection, $i\colon Y_0\to Y$ is the canonical immersion, and
$(g_0,g_0)\colon Y_0\to X\times_S X$ is the morphism
$\Delta_{X/S}g_0$ with components $g_0,g_0$:
$$
\xymatrix@C=2.2cm{ X\times_S X \ar[d]_{{\pr}_1} & Y_0
\ar[l]_-{h_0=(g_0,g_0)} \ar[d]^i \\ X & Y \ar[l]_-g
\ar@{-->}[lu]_-h }
$$
Since $h_0$ factors through the diagonal immersion $\Delta_{X/S}$ and
since $Y$ is in the first-order infinitesimal neighborhood of~$Y_0$
(\ie $\cal{J}^2=0$), the $h$ sought necessarily factor (uniquely)
through the first-order infinitesimal neighborhood of the diagonal. This
neighborhood is identified (as an $X$-prescheme by means of ${\pr}_1$)
with the spectrum $X'$ of the sheaf of algebras
$\cal{O}_X+\Omega^1_{X/S}$ (where the second term is regarded as a
square-zero ideal), the diagonal morphism $X\to X'$ corresponding to the
canonical augmentation of this sheaf of algebras. Put
$Y'=X'\times_X Y$, $Y'_0=Y'\times_Y Y_0=X'\times_X Y_0$, so that the
$h$ sought are in one-to-one correspondence with sections $u$ of~$Y'$
over~$Y$ extending a given section $u_0$ of~$Y'_0$ over~$Y_0$. One may
moreover identify $Y'$ with the spectrum of the sheaf of algebras on~$Y$
$\cal{A}=g^*(\cal{O}_X+\Omega^1_{X/S})=\cal{O}_Y+
g^*(\Omega^1_{X/S})$, and $Y'_0$ with the sheaf of algebras
$\cal{A}_0=\cal{A}\otimes_{\cal{O}_Y}\cal{O}_{Y_0}
=\cal{O}_{Y_0}+g_0^*(\Omega^1_{X/S})$. Then the section $u_0$ is
the one defined by the canonical augmentation of~$\cal{A}_0$ into~$\cal{O}_{Y_0}$.
Thus $\cal{P}(Y)$ identifies with the set of algebra homomorphisms
$\cal{A}\to\cal{O}_Y$ inducing the canonical augmentation
$\cal{A}_0\to\cal{O}_{Y_0}$. Now the algebra homomorphisms
$\cal{A}\to\cal{O}_Y$ correspond bijectively to homomorphisms of
modules $\cal{M}\to\cal{A}$ (putting, for simplicity,
$\cal{M}=g^*(\Omega^1_{X/S})$), and we are interested in those which
induce the {\emph{zero}} homomorphism $\cal{M}_0\to\cal{O}_{Y_0}$ (where
$\cal{M}_0=\cal{M}\otimes_{\cal{O}_Y}\cal{O}_{Y_0}$), \ie those which
map $\cal{M}$ into the ideal~$\cal{J}$ of the augmentation. We therefore
obtain the set
$\Hom_{\cal{O}_Y}(\cal{M},\cal{J})=\Hom_{\cal{O}_{Y_0}}(\cal{M}_0,\cal{J})$
(since $\cal{J}$ is annihilated by~$\cal{J}$). This is the desired
canonical bijection \eqref{eq:III.5.1.*}.

Taking into account the implication (i)$\To$(iii) in~\ref{III.3.1}, one
obtains:

\begin{corollary}
\label{III.5.2}
If $X$ is smooth over~$S$ (at least at the points of~$g_0(Y_0)$), then
$\cal{P}$ is even a \emph{principal homogeneous sheaf} under the
commutative sheaf of groups~$\cal{G}$, which in this case may also be
written
\begin{equation*}
\cal{G}=g_0^*(\goth{g}_{X/S})\otimes_{\cal{O}_{Y_0}}\cal{J}
\end{equation*}
where $\goth{g}_{X/S}$ is the sheaf on~$X$ dual to~$\Omega^1_{X/S}$,
\ie the \emph{tangent sheaf} (or \emph{sheaf of derivations}) of~$X$
relative to~$S$. (This last formula follows from the fact that
$\Omega^1_{X/S}$ is then locally free of finite type.)
\end{corollary}

In particular, this principal homogeneous sheaf determines a cohomology
class in~$\H^1(Y_0,\cal{G})$, whose vanishing is necessary and sufficient
for the existence of an $S$-morphism~$g$ extending~$g_0$. And if such an
extension exists, the set of all possible extensions is a homogeneous
space under the group~$\H^0(Y_0,\cal{G})$.

In applying the methods of formal geometry, the situation is most often
the following: two $S$-preschemes $X$ and~$Y$ are given, together with
a coherent ideal~$\cal{I}$ on~$S$; let $S_n$ be the closed sub-prescheme
of~$S$ defined by~$\cal{I}^{n+1}$, and put
\begin{equation*}
X_n=X\times_S S_n \quoi,\quad Y_n=Y\times_S S_n.
\end{equation*}
Suppose that one has an $S_n$-morphism
\begin{equation*}
g_n\colon Y_n\to X_n
\end{equation*}
(or, what amounts to the same thing, an $S$-morphism $Y_n\to X$, or
again an $S_{n+1}$-morphism $Y_n\to X_{n+1}$, since such a morphism
necessarily induces $Y_n\to X_n$), and one seeks to extend it to an
$S_{n+1}$-morphism
\begin{equation*}
g_{n+1}\colon Y_{n+1}\to X_{n+1}.
\end{equation*}
(If one can continue indefinitely, one thus obtains a morphism
$\widehat{Y}\to\widehat{X}$ for the formal preschemes completed from~$Y$
and~$X$ for the ideals $\cal{I}\cal{O}_Y$ and~$\cal{I}\cal{O}_X$.) One may
apply~\ref{III.5.1} after replacing $(S,X,Y,Y_0,g_0)$ by
$(S_{n+1},X_{n+1},Y_{n+1},Y_n,g_n)$. The sheaf~$\cal{G}$ becomes here
the sheaf of homomorphisms of modules from
$g_n^*(\Omega^1_{X_{n+1}/S_{n+1}})$ to
$\cal{J}=\cal{I}^{n+1}\cal{O}_Y/\cal{I}^{n+2}\cal{O}_Y$. Since~$\cal{J}$ is
annihilated by~$\cal{I}\cal{O}_Y$, one can replace
$g_n^*(\Omega^1_{X_{n+1}/S_{n+1}})$ by the sheaf it induces on~$Y_0$,
namely $h_0^*(\Omega^1_{X/S})$, where $h_0$ is the composite
$Y_0\to Y_n\to X_{n+1}$, or again the composite
$Y_0\to X_0\to X_{n+1}$, where $g_0\colon Y_0\to X_0$ is induced by~$g_n$.
Since the inverse image of~$\Omega^1_{X_{n+1}/S_{n+1}}$ on
$X_0=X_{n+1}\times_{S_{n+1}}S_0$ is~$\Omega^1_{X_0/S_0}$, one sees
that one also has
\begin{equation*}
\cal{G}=\SheafHom_{\cal{O}_{Y_0}} (g_0^*(\Omega^1_{X_0/S_0}),
\cal{I}^{n+1}\cal{O}_Y/\cal{I}^{n+2}\cal{O}_Y).
\end{equation*}
Thus one obtains the

\begin{corollary}
\label{III.5.3}
Let $S,X,Y,\cal{I},g_n$ be as above, and let $\cal{P}(g_n)$ be the sheaf
on~$Y$ whose sections on an open set~$U$ are the extensions $g_{n+1}$ of
$g_n$ to an $S_{n+1}$-morphism $Y_{n+1}\to X_{n+1}$. Then
$\cal{P}(g_n)$ is a formally principal homogeneous sheaf under the
sheaf of groups
\begin{equation*}
\cal{G}=\SheafHom_{\cal{O}_{Y_0}} (g_0^*(\Omega^1_{X_0/S_0}),
\gr_{\cal{I}\cal{O}_Y}^{n+1}(\cal{O}_Y)).
\end{equation*}
\end{corollary}

In particular:

\begin{corollary}
\label{III.5.4}
If in addition $X$ is smooth over~$S$ (at least at the points of~$g_0(Y_0)$),
then $\cal{P}(g_n)$ is even a principal homogeneous sheaf. In particular,
it defines an obstruction class in~$\H^1(Y_0,\cal{G})$, whose vanishing is
necessary and sufficient for the existence of a global extension~$g_{n+1}$
of~$g_n$. And if such an extension exists, the set of all global extensions
is a principal homogeneous space under~$\H^0(Y_0,\cal{G})$. Finally, in the
case under consideration, the sheaf~$\cal{G}$ may also be written
\begin{equation*}
\cal{G}=g_0^*(\goth{g}_{X_0/S_0})\otimes_{\cal{O}_{Y_0}}
\gr_{\cal{I}\cal{O}_Y}^{n+1}(\cal{O}_Y)
\end{equation*}
\end{corollary}

Proceeding step by step, one sees therefore that if all the
$\H^1(Y_0,\cal{G}_n)$ vanish (where
$\cal{G}_n=g_0^*(\goth{g}_{X_0/S_0})\otimes
\gr_{\cal{I}\cal{O}_Y}^{n}(\cal{O}_Y)$), then, starting with an arbitrary
$g_k$, one can extend it successively to $g_{k+1},\dots$. In particular,
if~$\cal{I}$ is nilpotent, one will be able to find an extension~$g$ of
$g_k$ to~$Y$. The condition that the $\H^1$ vanish is satisfied in
particular if~$Y_0$ is affine. Thus one obtains:

\begin{corollary}
\label{III.5.5}
In the statement of the theorem~\ref{III.3.1}, one obtains a necessary and
sufficient condition equivalent to the others by supposing that the~$Y'$
which occurs in \textup{(ii)} (or \textup{(ii~bis)}) is affine, and by
requiring the existence of a \emph{global extension}~$g$ of~$g_0$ to all
of~$Y'$.
\end{corollary}

It should be noted that the proof of~\ref{III.3.1} could not give us this
result directly.

An important case is that in which~$Y$ is \emph{flat} over~$S$; then one
has
\begin{equation*}
\gr^n(\cal{O}_Y)=\gr^n(\cal{O}_S)\otimes_{\cal{O}_{S_0}}\cal{O}_{Y_0}
\end{equation*}
and when, in addition, the $\gr^n(\cal{O}_S)$ are \emph{locally free on}~$S$,
one obtains
\begin{equation*}
\cal{G}_n=
\SheafHom_{\cal{O}_{Y_0}}(g_0^*(\Omega^1_{X_0/S_0}),\cal{O}_{Y_0})
\otimes_{\cal{O}_{S_0}}\gr^n(\cal{O}_S),
\end{equation*}
or again, if $\Omega^1_{X_0/S_0}$ is itself locally free (for example,
if $X$ is smooth over~$S$),
\begin{equation*}
\cal{G}_n=g_0^*(\goth{g}_{X_0/S_0})
\otimes_{\cal{O}_{S_0}}\gr^{n}(\cal{O}_S)
\end{equation*}
If, for example, $S$ is affine with affine ring~$A$, and $\cal{I}$ is
defined by an ideal~$I$ of~$A$, one obtains
\begin{equation*}
\H^i(Y_0,\cal{G}_n)=\H^i(Y_0,\cal{G}_0)\otimes_A\gr_I^n(A)
\end{equation*}
for every~$i$ (indeed, the question is local on~$S_0$, and one is reduced
to the case where one tensors by a free module). \emph{In this case, the
vanishing of~$\H^1(Y_0,\cal{G}_0)$ implies that all the obstructions to
the successive extensions of~$g_n$ vanish}. Thus one obtains:

\begin{corollary}
\label{III.5.6}
Let $(S,X,Y,\cal{I},g_n)$ be as above, suppose in addition that $X$ is
smooth over~$S$ and $Y$ is flat over~$S$, and finally that~$S$ is affine,
and that the $\gr^n(\cal{O}_S)=\cal{I}^n/\cal{I}^{n+1}$ are locally free.
Then the obstruction to constructing~$g_{n+1}$ lies in
$\H^1(Y_0,\cal{G}_0)\otimes_A\gr_I^{n+1}(A)$ (where~$A$ is the ring of~$S$,
$ I$ the ideal of~$A$ defining~$\cal{I}$), on putting
\begin{equation*}
\cal{G}_0=g_0^*(\goth{g}_{X_0/S_0}).
\end{equation*}
If $\H^1(Y_0,\cal{G}_0)=0$, then $g_n$ can be extended to an
$\widehat{S}$-morphism $\widehat{g}\colon\widehat{Y}\to\widehat{X}$.
\end{corollary}

Of course, this result would remain valid exactly as stated if, instead of
starting with ordinary $S$-preschemes $X$ and~$Y$, one started with formal
$\widehat{\cal{I}}$-adic $\widehat{S}$-preschemes $\goth{X}$ and~$\goth{Y}$.
It allows one to prove, for example, that certain formal schemes proper
over a complete local ring are in fact algebraic. Indeed, proceeding as
in the lemma~\ref{III.4.2}, one obtains:

\begin{corollary}
\label{III.5.7}
Under the conditions of~\ref{III.5.6}, if $g_0$ is an isomorphism, then so
is~$\widehat{g}$.
\end{corollary}

\noindent
(N.B. the same result holds for closed immersions.)

Thus one obtains:

\begin{proposition}
\label{III.5.8}
Let $A$ be a complete local ring with maximal ideal~$\goth{m}$ and residue
field~$k$, and let $\goth{X}$ and~$\goth{Y}$ be two $\goth{m}$-adic formal
preschemes over~$A$, flat over~$A$ (\ie for every~$n$, $X_n$ and $Y_n$ are
flat over~$A_n=A/\goth{m}^{n+1}$). Suppose that
$X_0=\goth{X}\otimes_A k$ is smooth over~$k$ and that
$\H^1(X_0,\goth{g}_{X_0/k})=0$. Then every $k$-isomorphism from~$Y_0$ onto~$X_0$
extends to an $A$-isomorphism from~$\goth{Y}$ onto~$\goth{X}$; this extension
is unique if, in addition, $\H^0(X_0,\goth{g}_{X_0/k})=0$.
\end{proposition}

This gives in particular a result on the \emph{uniqueness} of a formal
prescheme smooth over~$A$ reducing to a given prescheme~$X_0$ (provided
$\H^1(X_0,\goth{g}_{X_0/k})=0$). Moreover, if~$\goth{X}$ and~$\goth{Y}$
come from ordinary proper schemes over~$A$, say $X$ and~$Y$, then one knows,
by the existence theorem for sheaves in formal geometry (\cf the exposé
of the Bourbaki Seminar, no.~182\footnote{\Cf EGA III 5.4.1 for the proof.}),
that there is a one-to-one correspondence between the $A$-isomorphisms
$Y\isomto X$ and the $A$-isomorphisms of the formal completions, hence

\begin{corollary}
\label{III.5.9}
The preceding statement~\ref{III.5.8} remains valid on replacing
$\goth{X}$ and~$\goth{Y}$ by ordinary $A$-schemes $X$ and~$Y$, \emph{proper}
over~$A$.
\end{corollary}

Finally, when~$\goth{X}$ is a formal scheme proper over~$A$, and~$\goth{Y}$
is of the form~$\widehat{Y}$ where~$Y$ is an ordinary proper scheme over~$A$,
Proposition~\ref{III.5.8} gives sufficient conditions for finding an
isomorphism of~$\goth{X}$ with~$\widehat{Y}$, hence for the formal scheme
$\goth{X}$ to be in fact ``algebraic'' (\ie isomorphic to a~$\widehat{X}$,
$X$ an ordinary proper scheme over~$A$, which will then be canonically
determined). This is notably what happens if $X_0=\PP_k^r$ (or more
generally, if $X_0$ is a Severi--Brauer scheme, \ie becomes isomorphic to
the standard projective space over the algebraic closure of~$k$): every
formal scheme proper and flat over~$A$, with fiber~$\PP_k^r$, is
algebraizable, and more precisely is isomorphic to the $\goth{m}$-adic
formal completion of~$\PP_A^r$. In particular (thanks to the ``existence
theorem'') every ordinary proper scheme over~$A$, with fiber~$\PP_k^r$,
is isomorphic to~$\PP_A^r$ ($A$ being a complete local ring). Using
descent theory, one can prove that if~$A$ is not complete, $X$ becomes
isomorphic to~$\PP^r$ after making a finite étale extension
$A\to A'$ of the base (and in this form, the result remains valid for a
fiber which is a Severi--Brauer scheme).

\section{Global infinitesimal extension of smooth $S$-schemes}
\label{III.6}

Under the conditions of Theorem~\ref{III.4.1}, we propose to seek whether there exists a prescheme~$X$ smooth over~$Y$ such that
$X\times_YY_0$ is $Y_0$-isomorphic to~$X_0$, knowing that such a scheme ``exists locally on~$X_0$.'' Taking up again the method
of construction step by step, one is led to replace~$Y$ by the letter~$S$, to suppose that one is given a closed
sub-prescheme~$S_0$ of~$S$ defined by a sheaf of ideals~$\cal{I}$ (which it is no longer necessary to suppose locally
nilpotent), to introduce the closed sub-preschemes~$S_n$ of~$S$ defined by~$\cal{I}^{n+1}$, and to suppose that one is given
a sub-prescheme~$X_n$ smooth over~$S_n$. One proposes to find an $S_{n+1}$-prescheme~$X_{n+1}$ ``reducing to~$X_n$,'' i.e. equipped with an isomorphism
\begin{equation*}
X_{n+1}\times_{S_{n+1}}S_n\isomfrom X_n
\end{equation*}
which is \emph{smooth} over~$S_{n+1}$ (or, what amounts to the same thing by II~\SourceRef{II.2.1}{2.1}, \emph{flat} over~$S_{n+1}$). As we noted in~\ref{III.4}, such data amounts to giving a sheaf of algebras~$\cal{B}$ on $f^{-1}(\cal{O}_{S_{n+1}})$ (where $f$ is the underlying continuous map of the structural morphism $X_n\to S_n$), equipped with an augmentation $\cal{B}\to\cal{O}_{X_n}$ compatible with the augmentation $f^{-1}(\cal{O}_{S_{n+1}})\to f^{-1}(\cal{O}_{S_n})$, and satisfying two conditions (a) and (b) that we shall not rewrite, merely noting that they are of a \emph{local nature} on the topological space underlying~$X_n$. By~\ref{III.4.1} one knows that a solution exists locally. It is moreover unique up to (nonunique) isomorphism, at least locally. Let us begin by making this point precise:

\begin{proposition}
\label{III.6.1}
Let
$X_{n+1}$ over~$S_{n+1}$ reduce to~$X_n$ over~$S_n$. Then the sheaf (on the topological space underlying~$X_n$, or equivalently~$X_0$) of the $S_{n+1}$-automorphisms of~$X_{n+1}$ which induce the identity on~$X_n$ is canonically isomorphic to
\begin{equation*}
\cal{G}=\goth{g}_{X_0/S_0}\otimes_{\cal{O}_{S_0}}
gr_{\cal{I}}^{n+1}(\cal{O}_S)
\end{equation*}
(as a sheaf of groups).
\end{proposition}

Indeed, by~\ref{III.5.4} and~\ref{III.4.2} this sheaf is a principal homogeneous sheaf under~$\cal{G}$. Since it is equipped with a distinguished section (the identity automorphism of~$X_{n+1}$), it is therefore identified, as a sheaf of sets, with~$\cal{G}$. One must verify that this identification is compatible with the group structures. This is easy, and is moreover a special case of a more general result on the compatibility of the principal-bundle structures in~\ref{III.5.1} and~\ref{III.5.3} with composition of morphisms (a result that we do not state here, but which ought to appear in the hyperplodocus).

In particular, the sheaf on~$X_0$ of germs of automorphisms of~$X_{n+1}$ (with the structures just made explicit) is \emph{commutative}. It follows that if~$X_{n+1}'$ is another solution of the problem, isomorphic to~$X_{n+1}$ over the open~$U$ of~$X_0$, then the isomorphism from~$\SheafAut(X_{n+1})|U$ to~$\SheafAut(X_{n+1}')|U$ deduced by transport of structure from an isomorphism $X_{n+1}|U\isomto X_{n+1}'|U$ \emph{does not depend} on the choice of the latter. (It is in fact nothing other than the identity isomorphism of~$\cal{G}$, when the two automorphism sheaves are identified with~$\cal{G}$ by~\ref{III.6.1}.)

From~\ref{III.6.1} one deduces:

\begin{corollary}
\label{III.6.2}
Let~$X_{n+1}$ and~$X_{n+1}'$ be smooth over~$S_{n+1}$ and ``reducing to~$X_n$.'' Then the sheaf (on the space underlying~$X_0$) of $S_{n+1}$-isomorphisms from~$X_{n+1}$ to~$X_{n+1}'$ inducing the identity on~$X_n$ is naturally a principal homogeneous sheaf under~$\cal{G}$.
\end{corollary}

This indeed expresses that~$X_{n+1}$ and~$X_{n+1}'$ are locally isomorphic, and that the sheaf of germs of automorphisms of the first is~$\cal{G}$.

Now note that by~\ref{III.4.1} one can always find a covering~$(U_i)$ of~$X_n$ by open sets (which may be supposed affine), and for every~$i$ a smooth scheme~$X^i$ over~$S_{n+1}$ reducing to~$U_i$. Suppose for simplicity that~$X_n$ is \emph{separated}, so that the~$U_{ij}=U_i\cap U_j$ are still \emph{affine} open subsets of~$X_n$. Since the~$\H^1$ of such an open set, with values in the quasi-coherent sheaf~$\cal{G}$, is zero, it follows by \emph{the} corollary~\ref{III.6.2} that~$X^i|U_{ij}$ is isomorphic to~$X^j|U_{ij}$; let
$$
f_{ji}:X^i|U_{ij}\isomto X^j|U_{ij}
$$
be such an isomorphism. It is determined up to a section of~$\cal{G}$ on~$U_{ij}$. Put, for every triplet of indices,
$$
f_{ji}^{(k)}=f_{ji}|U_{ijk}\quad\text{where }U_{ijk}=U_i\cap U_j\cap U_k.
$$
If one had
\begin{equation*}
\label{eq:III.6.1}
\tag{1} {f_{kj}^{(i)}f_{ji}^{(k)}=f_{ki}^{(j)},}
\end{equation*}
it would follow that the~$X^i$ glue by the~$f_{ji}$, and hence define a solution~$X=X_{n+1}$ of the problem sought. Such a solution exists more generally if one can modify the~$f_{ji}$ into~$f_{ji}'$:
\begin{equation*}
\label{eq:III.6.2}
\tag{2} {f_{ji}'=f_{ji}g_{ji}\quad (g_{ji}\in\Gamma(U_{ij},\cal{G}))}
\end{equation*}
so that the~$f_{ji}'$ satisfy the transitivity condition above. This sufficient condition for the existence of a solution is also necessary, as one sees by recalling that such a solution~$X$ must, on each~$U_i$, be isomorphic to~$X^i$, which therefore allows one to choose isomorphisms
$$
f_i:X|U_i\isomto X^i
$$
and to define
$$
f_{ji}'=(f_j|U_{ij})(f_i|U_{ij})^{-1}:X^i|U_{ij}\isomto X^j|U_{ij}
$$
satisfying the gluing condition.

Now put
\begin{equation*}
\label{eq:III.6.3}
\tag{3}
{f_{ijk}=(f_{ki}^{(j)})^{-1}f_{kj}^{(i)}f_{ji}^{(k)},}
\end{equation*}
it is an automorphism of~$X^i|U_{ijk}$, which we identify with a section of~$\cal{G}$ by~\ref{III.6.1}. One observes that it is a $2$-\emph{cocycle} $f$ of the open covering~$\cal{U}=(U_i)$, with coefficients in~$\cal{G}$, by a small formal calculation left to the reader. The same calculation shows that, subject to~\eqref{eq:III.6.2}, the gluing condition~\eqref{eq:III.6.1} \emph{for the} $f_{ij}'$ is equivalent to the formula
\begin{equation*}
\label{eq:III.6.4}
\tag{4} {f=dg,}
\end{equation*}
where~$g=(g_{ij})$ is regarded as a $1$-cochain of~$\cal{U}$ with coefficients in~$\cal{G}$. Thus \emph{the necessary and sufficient condition for the existence of a solution of the problem is that the cohomology class in $\H^2(\cal{U},\cal{G})$ defined by the cocycle~\eqref{eq:III.6.3}} be zero. Note moreover that since~$\cal{U}=(U_i)$ is an affine covering of~$X_0$, which is a \emph{scheme}, $\H^2(\cal{U},\cal{G})$ identifies with $\H^2(X_0,\cal{G})$. It is immediate that the cohomology class thus obtained in~$\H^2(X_0,\cal{G})$ does not depend on the affine covering considered. We shall call it the \emph{obstruction class to extending~$X_n$ to a scheme~$X_{n+1}$ smooth over~$S_{n+1}$}.

Suppose this obstruction is zero. Then the argument sketched above shows that every solution~$X=X_{n+1}$ is isomorphic to a solution obtained by gluing from isomorphisms~$f_{ji}'$, which may be supposed to have the form~\eqref{eq:III.6.2}, the gluing condition being none other than~\eqref{eq:III.6.3}. The set of admissible~$g$ is therefore a principal homogeneous space under the group $\mathrm{Z}^1(\cal{U},\cal{G})$ of $1$-cocycles of~$\cal{U}$ with coefficients in~$\cal{G}$. Moreover, one sees at once that \emph{two cochains~$g$ and~$g'$} (such that $dg=dg'=f$) \emph{define isomorphic solutions if and only if the cocycle~$g-g'$ is of the form~$dh$}, where~$h=(h_i)\in\mathrm{C}^0(\cal{U},\cal{G})$. Thus one obtains:

\begin{theorem}
\label{III.6.3}
Let $(S,\cal{I},X_n)$ be as above, with~$X_n$ assumed separated\footnote{This condition is in fact unnecessary, and one can avoid the cocycle calculations above. \Cf the book by J. \textsc{Giraud}, \emph{Cohomologie Non Abélienne} (forthcoming from Springer Verlag, 1971). Compare Remarks~\ref{III.6.4}.}. Then one can define canonically an obstruction class in
$\H^2(X_0,\cal{G})$ (where~$\cal{G}$ is defined in~\ref{III.6.1}), whose vanishing is necessary and sufficient for the existence of a scheme~$X_{n+1}$ smooth over~$S_{n+1}$ reducing to~$X_n$. If this obstruction is zero, then the set of isomorphism classes, up to isomorphism (inducing the identity on~$X_n$), of $S_{n+1}$-preschemes~$X_{n+1}$ reducing to~$X_n$ is naturally a principal homogeneous space under~$\H^1(X_0,\cal{G})$.
\end{theorem}

\begin{remarks}
\label{III.6.4}
Starting from~\ref{III.6.1}, the arguments made here are purely formal, and are advantageously transcribed in the setting of local categories, or even of general fibered categories. The obstruction class to the existence of a ``global'' object of a category (in which one can find an object ``locally,'' and in which two objects are always ``locally isomorphic,'' the automorphism group of every object being commutative) obtained in this general context contains as a special case the ``second boundary homomorphism'' in an exact sequence of sheaves of not necessarily commutative groups, studied for example by Grothendieck in Kansas or T\^ohoku. The silly cocycle calculation made here should therefore be regarded as a makeshift, due to the absence of a satisfactory reference text.
\end{remarks}

\subsection{}
\label{III.6.5}
It should be noted that in~\ref{III.6.3} there is in general no distinguished element in the principal homogeneous space under~$\H^1(X_0,\cal{G})$ under consideration. This is reflected in particular by the fact that, after localizing on~$S$, one obtains a principal homogeneous sheaf on~$S_0$, with structural group~$\R^1f_*(\cal{G})$, which is not necessarily trivial, i.e. which defines a cohomology class in~$\H^1(S_0,\R^1f_*(\cal{G}))$ which is not necessarily zero. (This is when one supposes that the class~$d\in\H^2(X_0,\cal{G})$ is not zero, but is zero ``locally over~$S$,'' i.e. defines a zero section of~$\R^2f_*(\cal{G})$, i.e. a zero element in~$\H^0(S_0,\R^2f_*(\cal{G}))$.)

\subsection{}
\label{III.6.6}
For the moment one knows almost nothing about the general algebraic mechanism of the cohomology classes introduced in this number and their relations with the preceding number, and one knows nothing precise to say about them in the simplest particular cases, such as the case of abelian schemes over artinian rings\footnote{\label{III.6.6.p24}%
It is now known that this obstruction is always zero in this case [added in 2003 (MR): \cf F\ptbl Oort, \emph{Finite group schemes, local moduli for abelian varieties and lifting problems}, \emph{Algebraic Geometry Oslo 1970}, Wolters-Noordhoff, 1972, pp.\ptbl 223--254].}. One hopes that there will be people to work the question out thoroughly; it seems particularly interesting. It is intimately linked, in particular, to the ``module theory'' of algebraic structures.

\begin{corollary}
\label{III.6.7}
Suppose that $\H^2(X_0,\cal{G})=0$. Then an~$X_{n+1}$ exists, and it is unique up to isomorphism if moreover $\H^1(X_0,\cal{G})=0$.
\end{corollary}

In particular, one concludes, proceeding step by step (and observing that an affine scheme is acyclic for a quasi-coherent sheaf):

\begin{corollary}
\label{III.6.8}
Under the conditions of Theorem~\ref{III.4.1}, if~$X_0$ is affine, then there exists an~$X$ smooth over~$Y$ reducing to~$X_0$, and this~$X$ is unique up to (nonunique) isomorphism.
\end{corollary}

It should be noted that the direct proof of Theorem~\ref{III.4.1} could not have given this result.

\begin{corollary}
\label{III.6.9}
Under the conditions of~\ref{III.6.3}, suppose~$S$ affine with ring~$A$, $\cal{I}$ defined by an ideal~$I$ of~$A$, and finally that the $\gr^n_{\cal{I}}(\cal{O}_S)=\cal{I}^n/\cal{I}^{n+1}$ are locally free. Then $\H^i(X_0,\cal{G})$ identifies with
$$
\H^i(X_0,\cal{G}_0)\otimes_A\gr^{n+1}_I(A),
$$
where
$$
\cal{G}_0=\goth{g}_{X_0/S_0},
$$
so the obstruction class to extending~$X_n$ lies in $\H^2(X_0,\cal{G}_0)\otimes_A\gr^{n+1}_I(A)$, and, if it is zero, the set of isomorphism classes (up to isomorphism) of solutions is a principal homogeneous space under
$\H^1(X_0,\cal{G}_0)\otimes_A\gr^{n+1}_I(A)$.
\end{corollary}

In particular:

\begin{corollary}
\label{III.6.10}
Under the conditions of~\ref{III.6.9}, suppose
$$
\H^2(X_0,\goth{g}_{X_0/S_0})=0,
$$
then there exists an $\hat{I}$-adic formal scheme~$\goth{X}$ over the $\cal{I}$-adic formal completion~$\hat{S}$ of~$S$, which is ``smooth over~$S$'' (i.e. the~$\goth{X}_p$ are smooth over the~$S_p$) and which reduces to~$X_n$, i.e. is equipped with an isomorphism
$$
\goth{X}\times_S S_n\isomfrom X_n.
$$
If moreover $\H^1(X_0,\goth{g}_{X_0/S_0})=0$, then such a~$\goth{X}$ is unique up to isomorphism.
\end{corollary}

Indeed, one constructs $X_{n+1}$, $X_{n+2}$, and so on, step by step, whence~$\goth{X}$ by passing to the inductive limit of the~$X_i$. The assertion of uniqueness already appears in the preceding number.

\section[Application to the construction of formal schemes...]{Application to the construction of formal schemes
and smooth ordinary schemes over a complete local ring~$A$}
\label{III.7}
%
%
The results of the preceding~\No sometimes make it possible to
prove the existence
%
% original p. 83
%
of an $\goth{m}$-adic formal scheme over such a ring,
reducing to a given smooth scheme~$X_0$ over~$k$.
Let us distinguish two cases:

a) \emph{$A$ is ``of equal characteristic''} (this is the case
in particular if~$k$ has characteristic~$0$). Then one knows
that there exists a \emph{subfield of representatives of}~$A$,
\ie a subfield~$k'$ such that $A\to k$ induces an isomorphism~$k'\isomto
k$. \emph{Then there even exists a smooth ordinary scheme
over~$A$ reducing to~$X_0$}, namely~$X=X_0\otimes_kA$, with~$A$
regarded as an algebra over~$k$ by means of the homomorphism
$k\to k'\to A$ defined by~$k'$. It should nevertheless be noted
that this construction is not ``natural''; it is easy to convince
oneself (already in the case where~$A=k[t]/(t^2)$, the algebra of
dual numbers) that another lifting homomorphism~$k\to A$ (defined
in this instance by an absolute derivation of~$k$ into itself)
defines an~$X'$ over~$A$ which in general \emph{is not isomorphic
to~$X$} (if~$\H^1(X_0,\goth{g}_{X_0/k})\neq0$). It would moreover
be interesting to study (for~$k$ of characteristic~$0$, or
non-perfect and of characteristic~$p>0$) which smooth~$X$ over~$A$
are obtained in this way, and under what condition two homomorphisms
$k\to A$ define isomorphic $A$-schemes. Nevertheless, the existence
of~$k'$ suffices to entail that the first obstruction to lifting
$X_0$, which is
in~$\H^2(X_0,\goth{g}_{X_0/k})\otimes_k\goth{m}/\goth{m}^2$, is
necessarily zero. Of course, when one has then lifted~$X_0$ to~$X_1$
smooth over~$A/\goth{m}^2$, the new obstruction to constructing~$X_2$
will in general not be zero: it will be a function of a variable
element in a certain principal homogeneous space
under~$\H^1(X_0,\cal{G}_0)\otimes\goth{m}/\goth{m}^2$ and lies
in~$\H^2(X_0,\cal{G}_0)\otimes\goth{m}^2/\goth{m}^3$: it would
be appropriate to study the situation in detail\footnote{It is
no doubt described by the Kodaira--Spencer bracket operation (\cf
Cartan Seminar, 1960/61, Exp\ptbl 4).}.

b) \emph{$A$ is of unequal characteristics.} In this case, we
know nothing (except if by chance~$\H^2(X_0,\goth{g}_{X_0/k})=0$,
in which case one can construct a simple $\goth{m}$-adic formal
scheme over~$A$ reducing to~$k$). Even if~$A=\ZZ/p^2\ZZ$ and~$X_0$
is an ``abelian'' scheme of dimension~$2$, one does not know whether
it can be lifted to an~$X=X_1$ smooth over~$A$\footnote{This has
now been proved, \cf note~\ref{III.6.6.p24} page~\pageref{III.6.6.p24}.},
on the other hand, one has no example of an~$X_0$ which one has
proved does not come from an ordinary scheme~$X$ smooth over~$A$.
(I have the impression that such an example must exist, with~$X_0$
a projective surface)\footnote{\label{nIII.3}Such an example has since been
constructed by
J-P\ptbl Serre (Proc.\ Nat.\ Acad.\ Sc.\ USA \textbf{47} (1961),
\No 1, p\ptbl108--109), at least for certain
dimensions. D\ptbl Mumford gave an example (unpublished) with
an algebraic \emph{surface}.}. Let us simply point out that,
according to Cohen's theorem, there exists a Cohen $p$-ring~$B$
with residue field~$k$ and a homomorphism~$B\to A$ inducing the
identity isomorphism on the residue fields; consequently, the
``strongest'' lifting result would be obtained by taking~$A$
%
% original p. 84
%
a Cohen $p$-ring: if there exists a solution (ordinary or formal)
over such a ring, there exists one over every complete local ring
with residue field~$k$. In particular, since for a Cohen $p$-ring
$\goth{m}/\goth{m}^2$ is canonically identified with~$k$, one sees
that \emph{for every smooth scheme~$X_0$ over a field~$k$ of
characteristic~$p>0$, there exists a cohomology class
in~$\H^2(X_0,\goth{g}_{X_0/k})$} which is the first obstruction to
lifting~$X_0$ to a smooth scheme over a Cohen $p$-ring; one does
not know whether it can be nonzero\footnote{It can be nonzero,
as indicated in note~\ref{nIII.3} on page~\pageref{nIII.3}.}.

Even if one succeeds step by step in constructing the~$X_n$
reducing to~$X_0$, this generally gives only a \emph{formal}
scheme~$\goth{X}$ smooth over~$A$, reducing to~$X_0$. When~$X_0$
is proper over~$A$, the question remains whether~$\goth{X}$ is in
fact algebraizable, in order to obtain an \emph{ordinary} scheme
proper over~$A$ and simple over~$A$, reducing to~$X_0$. The only
known criterion (mentioned in the Bourbaki Seminar, and appearing
in the Elements, Chapter~III, 4.7.1) is the following: if~$\goth{X}$
is proper over~$A$, and if~$\cal{L}$ is an invertible sheaf on~$\goth{X}$
such that the sheaf induced~$\cal{L}_0$ on~$X_0$ is ample (\ie a suitable
tensor power~$\cal{L}_0^{\otimes n}$, $n>0$, comes from a projective
embedding of~$X_0$), then there exists a scheme~$X$ projective over~$A$,
and an ample invertible sheaf on~$X$, such that~$(\goth{X},\cal{L})$
is obtained from it by $\goth{m}$-adic completion. This therefore
leads us, given a locally free sheaf~$\cal{E}_0$ on~$X_0$ (which we
shall choose ample and invertible for our purpose), to extend it to
a locally free sheaf~$\cal{E}$ on~$\goth{X}$. For this, we are
reduced to constructing step by step locally free sheaves~$\cal{E}_n$
on the~$X_n$. The discussion is entirely analogous to that in
\No\ref{III.6}, (\cf remark~\ref{III.6.4}), the essential role being
played by the \emph{sheaf of automorphisms} of an~$\cal{E}_{n+1}$
which induce the identity on~$\cal{E}_n$: one immediately shows that
this sheaf is identified with
\begin{equation*}
\label{eq:III.7.*}\tag{$*$}{}
\cal{G}= \SheafHom_{\cal{O}_{X_0}}(\cal{E}_0,
\cal{E}_0\otimes\gr_{\cal{I}}^{n+1}(\cal{O}_X))=
\SheafHom_{\cal{O}_{X_0}}(\cal{E}_0,
\cal{E}_0)\otimes\gr_{\cal{I}}^{n+1}(\cal{O}_X)
\end{equation*}
which is again a sheaf of commutative groups. We find:

\begin{proposition}\label{III.7.1}
Let~$S$ be a prescheme endowed with a quasi-coherent sheaf of
ideals~$\cal{I}$, $X$ a prescheme over~$S$, $S_n$ the subprescheme
of~$S$ defined by~$\cal{I}^{n+1}$, and~$X_n=X\times_SS_n$ (for every
integer~$n$). Let~$\cal{E}_n$ be a locally free sheaf
%
% original p. 85
%
on~$X_n$; we propose to extend it to a locally free sheaf~$\cal{E}_{n+1}$
on~$X_{n+1}$. Then~$\cal{E}_n$ defines a canonical obstruction class
in~$\H^2(X_0,\cal{G})$, where~$\cal{G}$ is the quasi-coherent sheaf
given by the formula above, whose vanishing is necessary and
sufficient for the existence of an~$\cal{E}_{n+1}$ extending~$\cal{E}_n$.
If this class is zero, then the set of classes, up to isomorphism
(inducing the identity on~$\cal{E}_n$), of solutions~$\cal{E}_{n+1}$,
is a principal homogeneous space under~$\H^1(X_0,\cal{G})$.
\end{proposition}

This proposition gives rise to the usual corollaries. Let us point
out only that if~$X$ is \emph{flat} over~$S$, then one can write
$$
\cal{G}= \SheafHom_{\cal{O}_{X_0}}(\cal{E}_0, \cal{E}_0)
\otimes_{\cal{O}_{S_0}} \gr_{\cal{I}}^{n+1}(\cal{O}_S)
$$
whence, if~$S$ is affine of ring~$A$ and if the~$\cal{I}^n/\cal{I}^{n+1}$
are locally free, the sufficient condition
$$
\H^2(X_0,\cal{G}_0)=0 \qquad\text{with}\qquad
\cal{G}_0=\SheafHom_{\cal{O}_{X_0}}(\cal{E}_0, \cal{E}_0)
$$
for the existence of an~$\cal{E}_{n+1}$, and hence step by step for
the existence of successive extensions~$\cal{E}_m$ ($m=n,n+1,$
etc.).

Returning to the initial situation, we therefore find:

\begin{proposition}\label{III.7.2}
Let~$A$ be a complete local ring,~$\goth{X}$ a proper and flat formal
scheme over~$A$, such that~$X_0$ is projective and
$\H^2(X_0,\cal{O}_{X_0})=0$. Then there exists a scheme~$X$
projective over~$A$ whose $\goth{m}$-adic formal completion is
isomorphic to~$\goth{X}$.
\end{proposition}

Combining this with~\ref{III.6.10}, we find:

\begin{theorem}\label{III.7.3}
Let~$A$ be a complete local ring with residue field~$k$, $X_0$ a
projective and smooth scheme over~$k$, such that
$$
\H^2(X_0,\goth{g}_{X_0/k})=\H^2(X_0,\cal{O}_{X_0})=0
$$
Then there exists a smooth and projective scheme~$X$ over~$A$,
reducing to~$X_0$.
\end{theorem}

More generally, if one is given an~$X_n$ smooth over~$A_n=A/\goth{m}^{n+1}$
reducing to~$X_0$, then there exists an~$X$ smooth and proper over~$A$
and an isomorphism~$X\otimes_AA_n=X_n$.

\begin{corollary}\label{III.7.4}
Every smooth and proper curve over~$k$ is obtained by reduction from
a smooth and proper curve over~$A$.
\end{corollary}

It is this result which will be the essential tool (together with the
existence theorem for sheaves in Formal Geometry) for studying the
fundamental group of~$X_0$ by transcendental means.
%
% original p. 86
%

\end{document}
